% Flashcards template — preamble is fixed; replace placeholders and add cards.
\documentclass[avery5388,grid,frame]{flashcards}
\usepackage[utf8]{inputenc}

\usepackage{graphicx}
\usepackage{tabularx}
\usepackage{enumitem}
\usepackage{comment}
\usepackage{amsmath}

\cardfrontstyle[\large\slshape]{headings}
\cardbackstyle{empty}

\begin{document}
 \cardfrontfoot{<DECK TITLE>}

 % --- Style A: simple list ---------------------------------------------------
 \begin{flashcard}[<CATEGORY>]{<TERM>}
  \medskip
  \textbf{<LEAD-IN>:}
  \begin{itemize}
   \item <ITEM 1>
   \item <ITEM 2>
   \item <ITEM 3>
  \end{itemize}
 \end{flashcard}

 % --- Style B: centered prose, vertically filled -----------------------------
 \begin{flashcard}[<CATEGORY>]{<TERM>}\centering\vspace{\fill}
  <DEFINITION OR EXPLANATION>
  \vspace{\fill}
 \end{flashcard}

 % --- Style C: short statement + bullet list ---------------------------------
 \begin{flashcard}[<CATEGORY>]{<TERM>}\vspace{\fill}
  \begin{center}
    <SHORT STATEMENT OR DEFINITION>
  \end{center}
  \vspace{1em}
  \textbf{<LEAD-IN>:}
  \begin{itemize}
   \item <ITEM 1>
   \item <ITEM 2>
   \item <ITEM 3>
  \end{itemize}
  \vspace{\fill}
 \end{flashcard}

\end{document}
